% Drawing 74/07-11 -- folder 74, batch 07, archive page 133 (PDF page 134).
% Transcription: un hexagone avec un point central $a$, les arêtes marquées $\xi_1$, $\xi_1^*$ \uncertain{} ; un dessin entièrement raturé ; un hexagone contenant un losange centré ; un hexagone divisé par des diagonales ; quelques points et flèches.
% Kind: hexagons (centre a, centred rhombus, diagonals), struck drawing. Method 2 (schematic redraw from the facsimile). Status: partial.
% Notes: A hexagon with its central point a joined to the left and right vertices, its top side dashed, its other sides and the two spokes bold; the edge labels xi1 (left) and xi1* (right; reading uncertain, as in the transcription). A hexagon containing a centred rhombus: two opposite vertices of the rhombus are the left and right vertices of the hexagon, the other two are inner points, and its centre is marked. A hexagon divided by its three long diagonals, which meet at one computed point (the intersection of two of them; the third is drawn through it, as on the page). At the right, a few points, two of them joined, and a double arrow. Struck out and omitted: the drawing entirely scribbled over (below the first hexagon) and the letters scrawled under the first hexagon. Still differs: the edge labels of the first hexagon and the marks by them are a reading. Every edge is drawn between the same vertex points. Redrawn 2026-11-09 (issue 128): to check again.
% Coordinates: 1 unit = 1 mm = 1/100 of the crop width (bbox in page fractions: [0.6, 0.57, 0.9, 0.74]).
% Redrawn by Claude (claude-opus-5-5), issue #73, 2026-11-09; not checked by a human; not compiled (no TeX in the session).
% Rights: reproduces part of the work -- behind RENDER_TRANSCRIPTS, never published by itself.
\documentclass[tikz,border=3pt]{standalone}
\usepackage{amsmath,amssymb}
\usetikzlibrary{patterns,arrows.meta}
\tikzset{pencil/.style={black!45}, struck/.style={black!35}, mark/.style={circle,draw,inner sep=1.3pt}, pt/.style={circle,fill,inner sep=0.9pt}, opt/.style={circle,draw,fill=white,inner sep=1.1pt}}
\begin{document}
\begin{tikzpicture}[x=1mm,y=1mm,line cap=round,line join=round,font=\scriptsize]
\draw[dashed] (14.5,-5) -- (36,-7);
\draw[very thick] (14.5,-5) -- (14,-15);
\draw[very thick] (14,-15) -- (24,-21);
\draw[very thick] (24,-21) -- (35.5,-17);
\draw[very thick] (35.5,-17) -- (36,-7);
\draw[very thick] (27,-10) -- (14,-15);
\draw[very thick] (27,-10) -- (35.5,-17);
\node[pt] at (27,-10) {};
\node[inner sep=1pt] at (26,-7) {\small $a$};
\node[inner sep=1pt] at (7,-17) {\small $\xi_1$};
\node[inner sep=1pt] at (41,-14) {\small $\xi_1^*$};
\draw (55,-28) -- (64,-28);
\draw (64,-28) -- (67,-42);
\draw (67,-42) -- (63,-55);
\draw (63,-55) -- (53,-55);
\draw (53,-55) -- (50,-40);
\draw (50,-40) -- (55,-28);
\draw (50,-40) -- (59,-34);
\draw (59,-34) -- (67,-42);
\draw (67,-42) -- (58,-50);
\draw (58,-50) -- (50,-40);
\node[pt] at (59,-34) {};
\node[pt] at (58,-50) {};
\node[pt] at (50,-40) {};
\node[pt] at (67,-42) {};
\node[pt] at (59,-42) {};
\draw (53,-65) -- (62,-64);
\draw (62,-64) -- (68,-71);
\draw (68,-71) -- (62,-77);
\draw (62,-77) -- (52,-77);
\draw (52,-77) -- (48,-69);
\draw (48,-69) -- (53,-65);
\draw (57.1,-70.4) -- (53,-65);
\draw (57.1,-70.4) -- (62,-64);
\draw (57.1,-70.4) -- (68,-71);
\draw (57.1,-70.4) -- (62,-77);
\draw (57.1,-70.4) -- (52,-77);
\draw (57.1,-70.4) -- (48,-69);
\node[pt] at (57.1,-70.4) {};
\draw[-{Stealth[length=2mm]}] (69,-12) -- (80,-12);
\draw[-{Stealth[length=2mm]}] (69,-15) -- (79,-14.5);
\draw (85,-13) -- (91.5,-13);
\node[pt] at (85,-13) {};
\node[pt] at (91.5,-13) {};
\node[pt] at (84.5,-19) {};
\node[pt] at (91,-19.5) {};
\node[pt] at (85,-23.5) {};
\node[pt] at (91.5,-24) {};
\end{tikzpicture}
\end{document}
